% ============================================================
%  Adapting the Chatterbox Text-to-Speech Architecture to
%  Romanian via Phoneme-Level Preprocessing
%  IEEE Conference Format — IEEEtran.cls
% ============================================================
\documentclass[conference]{IEEEtran}

\usepackage{graphicx}
\usepackage{amsmath}
\usepackage{amssymb}
\usepackage{booktabs}
\usepackage{cite}
\usepackage{xcolor}
\usepackage{balance}
\usepackage[hidelinks]{hyperref}

% Figures live in ../resources/ relative to this file
\graphicspath{{../resources/}}

% ============================================================
\begin{document}

\title{Zero-Shot Romanian TTS via Phoneme-Level Fine-Tuning of Chatterbox Multilingual}

\author{\IEEEauthorblockN{Adrian-Bogdan STÂNEA}
\IEEEauthorblockA{
  \textit{Technical University of Cluj-Napoca, Romania}\\
  \{adriannstanea1@gmail.com\}}}

\maketitle

% ============================================================
\begin{abstract}
Romanian remains unsupported in most state-of-the-art (SOTA)
zero-shot text-to-speech (TTS) systems, including Chatterbox —
a 0.5-billion-parameter, MIT-licensed model based on a Llama
transformer backbone with conditional flow matching (CFM).
We present the first fine-tuning of the Chatterbox Multilingual
architecture for Romanian.
Our approach fine-tunes only the T3 text-to-token transformer on
the SWARA 1.0 corpus (21,304 utterances, $\sim$21 hours, 18 speakers),
while leaving the voice encoder and S3 speech generator frozen.
A central contribution is a phoneme-level text preprocessing
strategy that maps the two Romanian characters absent from the
existing 2,454-token grapheme vocabulary (\c{s}$\,{\to}\,$\textit{sh},
\c{t}$\,{\to}\,$\textit{ts}) to phonetically
equivalent existing tokens, fully preserving pretrained embedding
quality.
We demonstrate that the alternative approach — extending the
vocabulary with mean-initialized embeddings — induces posterior
collapse, a failure mode independently reported across multiple
languages and corroborated here experimentally.
After 13,000 training steps ($\sim$39 epochs) with a learning rate
of $5\times10^{-6}$, the model produces intelligible Romanian speech
from unseen reference speakers.
Systematic ablation over exaggeration and classifier-free guidance
(CFG) inference parameters on two zero-shot holdout speakers reveals
a consistent trade-off: low exaggeration favours intelligibility
(best CER $8.7\%$, WER $25.9\%$, UTMOS $3.46$) while higher
exaggeration improves speaker fidelity (cosine similarity up to
$0.875$).
These results establish phoneme-level preprocessing as a viable,
architecture-neutral strategy for adapting BPE-tokenised TTS models
to new languages without architectural modification.
\end{abstract}

\begin{IEEEkeywords}
text-to-speech synthesis, language adaptation, Romanian,
flow matching, zero-shot voice cloning, fine-tuning,
posterior collapse
\end{IEEEkeywords}

% ============================================================
\section{Introduction and Related Work}
\label{sec:intro}

\subsection{Introduction}

The past three years have fundamentally altered the landscape of
neural TTS.
VALL-E~\cite{valle} reframed speech synthesis as conditional language
modelling over discrete audio codec tokens, demonstrating zero-shot
voice cloning from three seconds of reference audio.
Concurrent and subsequent work explored flow matching as an
alternative generative backbone: rather than iteratively denoising
a signal, CFM learns a vector field that transports noise samples to
speech along a straight optimal-transport path defined by an ordinary
differential equation (ODE), enabling stable, high-fidelity synthesis
with far fewer inference steps~\cite{f5tts}.
CosyVoice~\cite{cosyvoice} combined supervised semantic tokens derived
from a multilingual ASR model with chunk-aware causal flow matching,
achieving human-parity naturalness across nine languages.

Despite this progress, Romanian remains a low-resource language in the
TTS domain.
Existing Romanian synthesis options span three regimes:
(i)~MMS-TTS-RON~\cite{mms}, a single-speaker VITS model derived from
religious-text data, covering the language but offering no voice
cloning capability and below-SOTA quality;
(ii)~FastPitch-RO~\cite{fastpitch_ro}, a non-autoregressive
duration-model system adapted to 18 speakers via TitaNet embeddings
on SWARA, offering multi-speaker synthesis but without zero-shot
generalisation;
(iii)~F5-TTS-RO~\cite{f5tts_ro}, the most recent and most directly
relevant work, which introduces a lightweight ConvNeXt input adapter
that keeps the F5-TTS backbone frozen, achieving WER $5.27\%$ and
speaker cosine similarity $0.90$ on Common Voice — though residual
English accent artefacts were noted in a subset of samples.

Chatterbox~\cite{chatterbox} is a family of SOTA open-source TTS
models (MIT license) that has not been adapted to Romanian.
Its original variant employs a 0.5B-parameter Llama backbone for
flow matching, trained on approximately 500,000 hours of cleaned
speech data, and provides zero-shot voice cloning with a unique
emotion exaggeration control parameter.
Notably, the Chatterbox architecture differs substantively from
F5-TTS: it uses a BPE grapheme tokeniser with 2,454 tokens (not
character padding), and processes speech through a two-stage T3/S3
pipeline rather than a single DiT backbone.
Adapting Chatterbox thus poses distinct challenges — and this work
constitutes the first published fine-tuning of the Chatterbox
Multilingual variant for Romanian.

Our contributions are:
\begin{enumerate}
\item A phoneme-level preprocessing strategy for adapting
  BPE-tokenised TTS models to new languages without vocabulary
  extension, demonstrated to prevent posterior collapse.
\item Fine-tuning of Chatterbox Multilingual on Romanian SWARA data,
  producing intelligible zero-shot synthesis.
\item Systematic ablation over exaggeration and CFG inference
  parameters, characterising the intelligibility--speaker-fidelity
  trade-off.
\end{enumerate}

\subsection{Related Work}

\subsubsection{Modern Zero-Shot TTS Architectures}

The neural codec language model paradigm, introduced by VALL-E~\cite{valle},
established that discrete audio tokens produced by a neural codec
(EnCodec) can serve as the prediction target for a language model
conditioned on a short acoustic prompt.
Seed-TTS~\cite{seedtts} extended this with reinforcement learning
for improved robustness and speaker similarity.
F5-TTS~\cite{f5tts} took a non-autoregressive approach: a Diffusion
Transformer (DiT) trained with flow matching, requiring no duration
model or explicit phoneme alignments, and demonstrating strong
code-switching on a 100k-hour multilingual corpus.
CosyVoice~\cite{cosyvoice} introduced supervised semantic tokens —
derived from a multilingual ASR encoder — as a bridge between text
and flow-matching-based speech generation, with its successor
CosyVoice~2 achieving human-parity streaming synthesis.

\subsubsection{Chatterbox Architecture Family}

Chatterbox~\cite{chatterbox} is built around a 0.5B-parameter Llama~3
transformer that maps text tokens to speech tokens (T3), followed by
a frozen S3 generator that synthesises mel-spectrograms via CFM.
All audio includes a PerTh perceptual watermark.
The Multilingual variant supports 23 languages (excluding Romanian)
with a 2,454-token grapheme tokeniser.
The Turbo variant (350M parameters, GPT-2 backbone, single-step
decoding) has documented instability on new languages due to its
BPE tokeniser and has been excluded from consideration here.

\subsubsection{Romanian TTS and Low-Resource Adaptation}

FastPitch-RO~\cite{fastpitch_ro} adapted the FastPitch architecture
to 18 Romanian speakers by conditioning on TitaNet-L speaker
embeddings extracted from SWARA.
F5-TTS-RO~\cite{f5tts_ro} demonstrated that a lightweight
learnable input adapter (ConvNeXt) could extend F5-TTS to Romanian
while keeping the entire model frozen, with a key finding that
full fine-tuning yielded better WER ($3.62\%$) but substantially
lower speaker similarity ($0.79$) compared to the adapter ($5.27\%$,
$0.90$), highlighting the inherent trade-off in low-resource adaptation.
MMS-TTS~\cite{mms} scales VITS-based synthesis to 1,100 languages
including Romanian, but operates as a single-speaker system.
Align2Speak~\cite{align2speak} proposed a GRPO-based framework that
uses ASR-derived reward signals to align multilingual TTS models to
low-resource languages, reporting CER reduction from $33\%$ to $3.94\%$
with only 30 minutes of Portuguese data.

The SWARA corpus~\cite{swara} provides 21 hours of high-quality
Romanian read speech from 17 speakers under CC-BY-SA 4.0, and is the
primary training resource for Romanian TTS research.

% ============================================================
\section{Theoretical Overview}
\label{sec:theory}

\subsection{Conditional Flow Matching for Speech Synthesis}

Flow matching~\cite{f5tts} defines a generative model by learning a
time-dependent vector field $v_t(x)$ that maps samples from a simple
prior distribution $p_0$ (typically Gaussian noise) to the data
distribution $p_1$ (speech mel-spectrograms) via an ODE:
\begin{equation}
  \frac{dx}{dt} = v_t(x), \quad x(0) \sim p_0.
\end{equation}
Training minimises the conditional flow matching (CFM) objective:
\begin{equation}
  \mathcal{L}_{\text{CFM}} =
    \mathbb{E}_{t,p_1(x_1),p_t(x|x_1)}
    \bigl\| v_t(x) - u_t(x \mid x_1) \bigr\|^2,
\end{equation}
where $u_t(x \mid x_1)$ is the conditional vector field for the
optimal-transport (OT) path between $x_0$ and $x_1$:
\begin{equation}
  u_t(x \mid x_1) = x_1 - x_0, \quad
  x_t = (1-t)\,x_0 + t\,x_1.
\end{equation}
Unlike diffusion models — which require dozens to hundreds of
denoising steps — OT-CFM interpolates linearly between noise and
data, enabling both faster training convergence and inference in as
few as 10--20 ODE steps with a numerical solver.
Classifier-free guidance (CFG) scales the conditional signal at
inference:
\begin{equation}
  \tilde{v}_t = v_t^{\varnothing} +
    \alpha \bigl(v_t^{\text{cond}} - v_t^{\varnothing}\bigr),
\end{equation}
where $\alpha$ is the CFG scale and $v_t^{\varnothing}$ is the
unconditional vector field.
Higher $\alpha$ strengthens the conditioning signal at the cost of
increased mode-seeking behaviour.

\subsection{The Chatterbox Architecture}

Chatterbox processes text-to-speech synthesis in four stages,
summarised as follows.

\textbf{Stage 1 — EnTokenizer.}
Input text is converted to a sequence of integer token IDs by a
BPE grapheme tokeniser containing 2,454 tokens covering Latin,
Cyrillic, and other scripts across 23 languages.
Text is lowercased and punctuation-normalised before tokenisation.

\textbf{Stage 2 — T3 (Text-to-Token Transformer).}
The T3 model is a 500M-parameter Llama~3-based autoregressive
transformer conditioned on (a) the text token sequence and (b) a
speaker embedding extracted by the Voice Encoder from a short
reference audio clip.
It predicts a sequence of discrete speech tokens representing the
phonetic content and prosody of the target utterance.
This is the \emph{only trainable component} during fine-tuning; all
other modules remain frozen.

\textbf{Stage 3 — S3Gen (Speech-to-Signal Generator).}
S3Gen synthesises a mel-spectrogram from the speech token sequence
via CFM, additionally conditioned on the prompt mel-spectrogram and
a speaker embedding.
The CFM backbone is an OT-flow-matching model with a Llama
transformer as the vector-field estimator.
Because S3Gen is frozen, it retains its full pretrained capacity for
high-fidelity acoustic generation.

\textbf{Stage 4 — Vocoder.}
A neural vocoder (also frozen) converts the mel-spectrogram to a
time-domain waveform at 24 kHz.

\textbf{Inference parameters.}
Two parameters are exposed at inference time:
\begin{itemize}
  \item \emph{Exaggeration} (exag $\in [0,1]$): scales the emotional
    intensity of generated speech by modulating the conditioning
    signal strength in the T3 model.
  \item \emph{Classifier-free guidance scale} (cfg $\in \mathbb{R}$):
    the $\alpha$ parameter in Eq.~(4), controlling the trade-off
    between text-fidelity and acoustic diversity.
\end{itemize}

\subsection{Evaluation Metrics}

\textbf{Intelligibility — CER and WER.}
Synthesised audio is transcribed using OpenAI Whisper
medium~\cite{whisper}, configured with \texttt{language=``ro''},
\texttt{beam\_size=5}, and \texttt{best\_of=5}; the transcript is
compared against the reference text with standard
dynamic-programming alignment.
CER operates at the character level; WER at the word level.
Both metrics are lower-is-better.

\textbf{Speech quality — UTMOS.}
The UTokyo-SaruLab MOS predictor~\cite{utmos} is a neural network
trained on the VoiceMOS Challenge 2022 dataset to predict Mean
Opinion Score (MOS, 1--5 scale) from a waveform without human
listeners.
UTMOS generalises reasonably across languages and recording
conditions, making it suitable as an automatic quality proxy.

\textbf{Speaker similarity.}
Speaker embeddings are extracted from both the synthesised audio
and the reference speaker clip using the same encoder used by
Chatterbox's Voice Encoder.
Cosine similarity between the two embedding vectors quantifies
voice-cloning fidelity, ranging from 0 (orthogonal) to 1 (identical).

% ============================================================
\section{Methodology}
\label{sec:method}

\subsection{Dataset: SWARA 1.0}

The SWARA Speech Corpus~\cite{swara} is the primary open resource
for Romanian TTS research.
We use SWARA 1.0, comprising 21,304 utterances recorded from 43
speaker-session variants (18 unique speakers across multiple
recording rounds) at 22,050\,Hz.
Utterances contain an average of 10.4 words (62.4 characters), with
64 unique characters including all five Romanian diacritic classes.
Audio files were converted from the original format to the LJSpeech
directory layout (WAV + CSV metadata) for compatibility with the
training pipeline.

Dataset splits follow a zero-shot evaluation design: two speakers
(BAS and SGS, spanning three recording rounds each — 5 speaker-session
variants totalling 2,484 utterances) are held out entirely for
evaluation.
The remaining utterances are split 90/10 per speaker into
training (18,740 samples, 38 speaker-variants) and validation
(80 samples, 36 speaker-variants).
This strategy ensures that both the acoustic characteristics and
the recording-condition variation of the holdout speakers are
completely unseen during training, providing a rigorous zero-shot
voice-cloning test.

Romanian presents specific phonological challenges for adaptation.
The language uses five diacritic characters: \u{a} (central
mid-vowel, similar to English schwa), \^{a} and \^{\i} (close central
unrounded vowel, a phone absent from most Western European languages),
\c{s} (palato-alveolar fricative, as in English ``shoe''), and
\c{t} (affricate /ts/, as in English ``cats'').
Two of these (\c{s}, \c{t}) are absent from the 23-language
grapheme vocabulary, motivating our preprocessing strategy.
Contemporary Romanian orthography admits both cedilla variants
(\c{s}, \c{t}) and comma-below variants for the same phones; our
preprocessor handles both Unicode forms to ensure full corpus coverage.

\subsection{Language Adaptation: Three Iterations}

\subsubsection{Iteration 1 — Vocabulary Extension (Failed)}

The community fine-tuning toolkit for Chatterbox recommends
extending the tokeniser vocabulary with missing characters for new
languages.
Following this approach, we added five Romanian tokens
(\c{s}, \c{t}, \c{S}, \c{T}, \u{A}),
expanding the vocabulary from 2,454 to 2,459 tokens.
The T3 model's embedding layer and output projection were resized;
new token embeddings were initialised with the mean of all existing
embeddings (mean initialisation).

Fine-tuning with this configuration caused the model to produce
unintelligible speech within a few hundred training steps.
Analysis indicates \emph{posterior collapse}: the T3 model ignores
the text conditioning path in favour of the stronger audio signal
from the voice encoder.
Mean-initialised embeddings are non-discriminative — positioned at
the centroid of the embedding space, equidistant from all existing
token clusters — providing a far weaker conditioning signal than the
well-trained audio encoder.
The model exploits this imbalance and degenerates into an
autoencoder that copies the reference audio regardless of text input.
This failure was independently reported by community users of the
same toolkit for Turkish, Norwegian, German, and Arabic, and
confirmed as a structural problem by the toolkit maintainer.

\subsubsection{Iteration 2 — Phoneme Mapping with Unstable Hyperparameters}

Two key observations motivated an alternative strategy:
(i)~only two Romanian characters (\c{s}, \c{t}) are
genuinely absent from the 2,454-token vocabulary — the vowels
\u{a}, \^{a}, \^{\i} are already present;
(ii)~both missing characters have phonetically exact equivalents
already in the vocabulary: \c{s}$\,{\leftrightarrow}\,$\textit{sh}
(IPA /\textrm{S}/) and \c{t}$\,{\leftrightarrow}\,$\textit{ts}
(IPA /\textrm{ts}/).
We therefore implemented a phoneme-level preprocessor that replaces
the two absent consonants with their existing multi-character
equivalents before tokenisation, while leaving all other characters
unchanged.
This preserves the full 2,454-token vocabulary without any
mean-initialised embeddings, retaining the discriminative quality of
all pretrained representations.

However, the initial training run at the same learning rate used in
Iteration~1 ($1\times10^{-5}$, batch size~32) produced unstable
outputs — the model generated unintelligible speech after a limited
number of training steps, indicating that the T3 transformer is
sensitive to learning-rate magnitude during fine-tuning, independently
of the vocabulary configuration.

\subsubsection{Iteration 3 — Stabilised Phoneme Mapping (Final Configuration)}

The phoneme-level preprocessing from Iteration 2 was retained; the
primary changes were a halving of the learning rate from
$1\times10^{-5}$ to $5\times10^{-6}$ and a reduction of the batch
size from 32 to 16.
Table~\ref{tab:iterations} summarises the three iterations.

\begin{table}[h]
\centering
\caption{Comparison of the three adaptation iterations.}
\label{tab:iterations}
\begin{tabular}{lccc}
\toprule
& \textbf{Iter.~1} & \textbf{Iter.~2} & \textbf{Iter.~3} \\
\midrule
Vocabulary size & 2,459 & 2,454 & 2,454 \\
Preprocessing & None & Phoneme & Phoneme \\
Embedding init & Mean & Pretrained & Pretrained \\
Learning rate & $1\times10^{-5}$ & $1\times10^{-5}$ & $5\times10^{-6}$ \\
Batch size & 32 & 32 & 16 \\
Outcome & Collapse & Unstable & Stable \\
\bottomrule
\end{tabular}
\end{table}

The phoneme mappings applied in Iterations 2 and 3 are summarised
in Table~\ref{tab:phoneme_map}.

\begin{table}[h]
\centering
\caption{Romanian phoneme preprocessing mappings.}
\label{tab:phoneme_map}
\begin{tabular}{lccc}
\toprule
Character & IPA & Mapped to & Token ID \\
\midrule
\c{s}, \c{S} & /\textrm{sh}/ & \textit{sh} & 120 \\
\c{t}, \c{T} & /\textrm{ts}/ & \textit{ts} & 192 \\
\u{a} & /$\partial$/ & \textit{kept} & 2413 \\
\^{a} & /\textrm{1}/ & \textit{kept} & 395 \\
\^{\i} & /\textrm{1}/ & \textit{kept} & 407 \\
\bottomrule
\end{tabular}
\end{table}

\subsection{Training Setup}

Only the T3 transformer is fine-tuned; the Voice Encoder and S3Gen
remain frozen throughout.
Training was conducted on an NVIDIA V100 (32\,GB VRAM) using FP16
precision.
Key hyperparameters: learning rate $5\times10^{-6}$, batch size~16,
gradient accumulation~1, maximum text length~256 tokens, maximum
speech length~850 tokens, prompt duration~3.0\,s.
We trained for 13,000 global steps ($\approx$39 epochs), saving
checkpoints every 100 steps (retaining the last 10).
Training was monitored via TensorBoard; no learning-rate schedule
was applied.

\subsection{Evaluation Protocol}

Each of the two holdout speakers (BAS, SGS) was evaluated under
four parameter combinations: exaggeration $\in \{0.3, 0.7\}$
$\times$ CFG scale $\in \{0.3, 0.5\}$, yielding eight experimental
conditions.
For each condition, 50 test utterances were synthesised using a
3-second speaker reference clip, producing 400 total samples.
CER and WER were computed via transcription with OpenAI Whisper
medium~\cite{whisper} (\texttt{language=``ro''}, \texttt{beam\_size=5},
\texttt{best\_of=5}).
UTMOS~\cite{utmos} and speaker cosine similarity were computed
per sample.
All statistics are reported as mean $\pm$ standard deviation over
the 50-sample per-condition pool.

% ============================================================
\section{Evaluation and Experimental Results}
\label{sec:results}

\subsection{Training Dynamics}

\begin{figure}[t]
\centering
\includegraphics[width=\columnwidth]{training_dynamics/loss_curves.png}
\caption{Training loss over 13,000 steps ($\approx$39 epochs).
  The loss decreases monotonically from 5.98 to 0.11, with rapid
  initial descent and gradual plateau after step $\sim$10,000.}
\label{fig:loss}
\end{figure}

Fig.~\ref{fig:loss} shows the training loss curve.
The model converges monotonically, with an initial sharp descent
from 5.98 to approximately 1.6 within the first 4,000 steps,
followed by a gradual refinement phase that reaches a minimum of
0.11 at step 13,000 — a 98\% total reduction.
The smooth trajectory, with no oscillation or divergence, confirms
that the learning-rate reduction implemented in Iteration~3 provides
stable training dynamics.
The plateau observed after step $\sim$10,000 suggests the model
has captured the core Romanian phonetic patterns within the
fixed-vocabulary representation; extended training at a lower
learning rate remains a promising avenue for further refinement.

\subsection{Intelligibility}

\begin{figure}[t]
\centering
\includegraphics[width=\columnwidth]{ablation_analysis/ablation_v1_cer_wer.png}
\caption{CER (left) and WER (right) across all eight experimental
  conditions. BAS speaker (blue) consistently achieves lower error
  rates than SGS (orange). Low exaggeration (Exag=0.3) generally
  favours intelligibility.}
\label{fig:cer_wer}
\end{figure}

Table~\ref{tab:metrics} and Fig.~\ref{fig:cer_wer} present
intelligibility results across all conditions.
The best intelligibility is achieved by BAS at exag=0.3, cfg=0.5
(CER $8.7\%$, WER $25.9\%$).
BAS consistently outperforms SGS across all parameter combinations:
aggregated mean CER $11.9\%$ vs $15.0\%$, mean WER $31.2\%$ vs
$38.9\%$ (Table~\ref{tab:speaker_comparison}).

The effect of the exaggeration parameter is clear for BAS: low
exaggeration (0.3) yields CER $\approx9$--$10\%$, while high
exaggeration (0.7) raises CER to $\approx10$--$19\%$.
The anomalously high variance at BAS exag=0.7, cfg=0.5
(CER std $0.49$) indicates occasional catastrophic generation
failures under this combination, suggesting that high-exaggeration
generation in BAS is more susceptible to inference instability.
For SGS, the exaggeration effect on intelligibility is more modest,
with CER ranging from $11.7\%$ to $17.3\%$ across conditions.

CFG scale has a secondary but observable effect: for BAS at low
exaggeration, cfg=0.5 reduces both CER and WER relative to cfg=0.3.
This is consistent with the theoretical expectation that stronger
CFG reinforces the text-conditioning signal, improving phoneme
reproduction accuracy.

% table_full_metrics spans both columns due to width
\begin{table*}[t]
\centering
\caption{Objective evaluation metrics for Romanian TTS (mean $\pm$ std).
  Best intelligibility per metric is in \textbf{bold}.}
\label{tab:metrics}
\begin{tabular}{lcccc}
\toprule
Condition & CER$\downarrow$ & WER$\downarrow$ & UTMOS$\uparrow$ & Spk.~Sim.$\uparrow$ \\
\midrule
BAS (exag=0.3, cfg=0.3) & $0.101 \pm 0.143$ & $0.279 \pm 0.233$ & $3.393 \pm 0.570$ & $0.810 \pm 0.075$ \\
BAS (exag=0.3, cfg=0.5) & $\mathbf{0.087 \pm 0.141}$ & $\mathbf{0.259 \pm 0.222}$ & $\mathbf{3.460 \pm 0.306}$ & $0.816 \pm 0.058$ \\
BAS (exag=0.7, cfg=0.3) & $0.100 \pm 0.116$ & $0.312 \pm 0.254$ & $3.123 \pm 0.347$ & $0.824 \pm 0.054$ \\
BAS (exag=0.7, cfg=0.5) & $0.188 \pm 0.490$ & $0.399 \pm 0.484$ & $3.120 \pm 0.326$ & $0.830 \pm 0.050$ \\
SGS (exag=0.3, cfg=0.3) & $0.165 \pm 0.233$ & $0.373 \pm 0.297$ & $2.975 \pm 0.504$ & $0.830 \pm 0.148$ \\
SGS (exag=0.3, cfg=0.5) & $0.147 \pm 0.153$ & $0.394 \pm 0.256$ & $3.082 \pm 0.302$ & $0.863 \pm 0.084$ \\
SGS (exag=0.7, cfg=0.3) & $0.117 \pm 0.112$ & $0.372 \pm 0.291$ & $3.105 \pm 0.233$ & $\mathbf{0.875 \pm 0.040}$ \\
SGS (exag=0.7, cfg=0.5) & $0.173 \pm 0.193$ & $0.417 \pm 0.287$ & $3.094 \pm 0.216$ & $0.872 \pm 0.053$ \\
\bottomrule
\end{tabular}
\end{table*}

\subsection{Speech Quality}

\begin{figure}[t]
\centering
\includegraphics[width=\columnwidth]{ablation_analysis/abalation_utmos.png}
\caption{UTMOS quality scores for BAS (left) and SGS (right) across
  exaggeration and CFG parameter combinations. BAS at low
  exaggeration achieves the highest perceived quality (3.39--3.46).}
\label{fig:utmos}
\end{figure}

Fig.~\ref{fig:utmos} shows UTMOS quality scores as a heatmap over
the parameter grid for each speaker.
The best quality is achieved by BAS at exag=0.3, cfg=0.5
(UTMOS $3.46$).
Low exaggeration consistently yields higher UTMOS for BAS
($3.39$--$3.46$ vs $3.12$ at exag=0.7), while the effect is
less pronounced for SGS ($2.98$--$3.08$ vs $3.09$--$3.11$).
BAS achieves higher quality overall (aggregated mean $3.27$
vs $3.06$ for SGS).

UTMOS scores in the $3.0$--$3.5$ range correspond to ``fair to
good'' perceived quality on the 1--5 scale — confirming intelligible
synthesis while indicating headroom for improvement relative to the
Chatterbox base model's performance on English.
The quality advantage of low exaggeration aligns with the
intelligibility findings, suggesting that exaggeration uniformly
degrades both dimensions of perceived output quality.

\subsection{Speaker Fidelity}

\begin{figure}[t]
\centering
\includegraphics[width=\columnwidth]{ablation_analysis/ablation_v1_speaker_similarity.png}
\caption{Speaker cosine similarity across conditions. SGS (orange)
  consistently achieves higher fidelity than BAS (blue). Higher
  exaggeration improves speaker preservation for both speakers.}
\label{fig:speaker_sim}
\end{figure}

Fig.~\ref{fig:speaker_sim} reveals a pattern complementary to the
intelligibility results.
SGS consistently achieves higher speaker cosine similarity
($0.830$--$0.875$) than BAS ($0.810$--$0.830$), inverting the
intelligibility ranking.
Critically, higher exaggeration improves speaker similarity for both
speakers: BAS increases from $\approx0.813$ (exag=0.3) to
$\approx0.827$ (exag=0.7); SGS from $\approx0.847$ to
$\approx0.873$.
This creates a clear parameter-driven trade-off: the settings that
maximise intelligibility degrade speaker fidelity, and vice versa.

The positive effect of CFG on speaker similarity is also observable:
higher cfg=0.5 slightly outperforms cfg=0.3 at all conditions for
SGS, while the effect for BAS is mixed.
All eight conditions exceed cosine similarity $0.80$, confirming
that the frozen Voice Encoder successfully transfers speaker
characteristics to unseen Romanian speakers via the fine-tuned T3.

\subsection{Speaker-Level Comparison}

\begin{table}[h]
\centering
\footnotesize
\setlength{\tabcolsep}{4pt}
\caption{Speaker-level metrics (mean $\pm$ std, across all conditions).}
\label{tab:speaker_comparison}
\begin{tabular}{lcccc}
\toprule
Spk. & CER & WER & UTMOS & Sim. \\
\midrule
BAS & $0.119{\pm}0.040$ & $0.312{\pm}0.054$ & $3.274{\pm}0.154$ & $0.820{\pm}0.007$ \\
SGS & $0.150{\pm}0.021$ & $0.389{\pm}0.018$ & $3.064{\pm}0.052$ & $0.860{\pm}0.018$ \\
\bottomrule
\end{tabular}
\end{table}

Table~\ref{tab:speaker_comparison} aggregates results across all
parameter conditions per speaker.
BAS exhibits stronger intelligibility and quality (CER $\pm$ std:
$0.119 \pm 0.040$; UTMOS $3.274 \pm 0.154$), while SGS demonstrates
stronger speaker preservation ($0.860 \pm 0.018$ vs $0.820 \pm 0.007$).
The standard deviations indicate that SGS is more consistent across
parameter conditions (lower CER and WER std) while BAS is more
sensitive to parameter choices — consistent with the high-variance
failure mode observed at BAS exag=0.7, cfg=0.5.

These complementary speaker profiles suggest that the fine-tuned T3
model generalises differently to different voice types: BAS produces
more legible text reproduction, while SGS more faithfully preserves
voice identity.
This speaker-dependent behaviour may reflect differences in the
speakers' acoustic characteristics relative to the training set
speakers, or in the saliency of their voice embeddings to the frozen
S3Gen generator.

\subsection{Parameter Interaction Analysis}

The ablation results across the exaggeration and CFG axes reveal
structured main effects that provide guidance for practitioners.
Table~\ref{tab:maineffects} summarises the mean performance for
each parameter level, averaged across the other dimension and across
both speakers.

\begin{table}[h]
\centering
\caption{Main effects of exaggeration and CFG on each metric
  (averaged over the other factor and both speakers).}
\label{tab:maineffects}
\begin{tabular}{llcccc}
\toprule
Factor & Level & CER$\downarrow$ & WER$\downarrow$ & UTMOS$\uparrow$ & Sim.$\uparrow$ \\
\midrule
Exag & 0.3 & 0.125 & 0.326 & 3.228 & 0.830 \\
Exag & 0.7 & 0.145 & 0.377 & 3.110 & 0.850 \\
\midrule
CFG  & 0.3 & 0.120 & 0.334 & 3.149 & 0.835 \\
CFG  & 0.5 & 0.149 & 0.368 & 3.189 & 0.845 \\
\bottomrule
\end{tabular}
\end{table}

The exaggeration parameter has the strongest effect on intelligibility:
increasing exaggeration from 0.3 to 0.7 raises mean CER by $\approx2.0$ pp
and WER by $\approx5.1$ pp, while reducing UTMOS by $\approx0.12$ points.
In contrast, speaker similarity \emph{increases} with exaggeration
($+0.02$), indicating the parameter simultaneously trades intelligibility
for speaker fidelity.
The CFG parameter exhibits a qualitatively similar but weaker pattern:
higher CFG marginally degrades CER/WER while improving speaker
similarity.
Critically, the two parameters are not fully collinear in their
effects — the interaction between high exaggeration and high CFG is
associated with the anomalously high CER variance observed for BAS
(std $0.49$), suggesting a synergistic instability under combined
strong conditioning.

For Romanian TTS applications prioritising intelligibility (e.g.,
voice assistants, accessibility tools), we recommend exag=0.3,
cfg=0.5, which achieves the best CER and WER.
For applications prioritising voice identity fidelity (e.g., voice
cloning, character synthesis), exag=0.7, cfg=0.3 or cfg=0.5 are
preferable, with cfg=0.3 providing better stability.

\subsection{Comparison with Prior Work}

Table~\ref{tab:priorwork} contextualises our results against the
two closest comparable systems.

\begin{table}[h]
\centering
\caption{Comparison with Romanian TTS prior work.
  (--) denotes capability not supported.}
\label{tab:priorwork}
\begin{tabular}{lccc}
\toprule
System & WER~(\%)$\downarrow$ & Spk.~Sim.$\uparrow$ & Zero-shot \\
\midrule
MMS-TTS-RON~\cite{mms}            & 5.77  & --    & No  \\
F5-TTS-RO (adapter)~\cite{f5tts_ro} & 5.27  & 0.90  & Yes \\
F5-TTS-RO (full FT)~\cite{f5tts_ro} & 3.62  & 0.79  & Yes \\
\midrule
\textbf{Ours (BAS, best)} & \textbf{25.9} & \textbf{0.816} & \textbf{Yes} \\
\textbf{Ours (SGS, best)} & \textbf{37.2} & \textbf{0.875} & \textbf{Yes} \\
\bottomrule
\end{tabular}
\end{table}

Our WER is higher than both prior systems, but direct comparison
requires caution.
F5-TTS-RO was evaluated on Common Voice Romanian, which comprises
shorter, crowdsourced utterances; our evaluation uses SWARA holdout
utterances of greater length and linguistic complexity.
Furthermore, different ASR transcription systems were used.
Speaker similarity is broadly comparable: our best condition
achieves $0.875$ vs the F5-TTS-RO adapter's $0.90$, and substantially
exceeds F5-TTS-RO full fine-tuning ($0.79$).

The key architectural distinction is that our approach requires
\emph{no model changes whatsoever}: no adapters, no new modules,
no vocabulary extension.
Only a 30-line text preprocessor and conservative fine-tuning of
the existing T3 component are needed — a substantially lower
implementation barrier than adapter-based methods.
The higher WER also reflects limited training scale ($\sim$21 hours,
39 epochs) and suggests that further training would close the gap
with prior work.

% ============================================================
\section{Conclusions and Future Work}
\label{sec:conclusion}

We have presented the first fine-tuning of the Chatterbox Multilingual
architecture for Romanian, demonstrating that phoneme-level text
preprocessing enables stable fine-tuning of a BPE-tokenised
Llama-based flow-matching model without any modification to the model
architecture or vocabulary.
The central finding is that mean-initialised vocabulary extension
causes posterior collapse — a failure that is language-agnostic and
reproducible — while phoneme mapping to existing tokens fully avoids
this problem by preserving the discriminative quality of pretrained
embeddings.
Systematic ablation over exaggeration and CFG inference parameters
reveals a consistent intelligibility--speaker-fidelity trade-off
that practitioners can exploit according to their application
requirements: low exaggeration (0.3) and moderate CFG (0.5) maximise
intelligibility (CER $8.7\%$, WER $25.9\%$), while higher
exaggeration (0.7) improves voice-cloning fidelity (cosine
similarity $\geq 0.87$).
All evaluated configurations produce intelligible Romanian speech
with speaker similarity above $0.80$ for unseen holdout speakers.

Several directions remain for future work.
The training loss plateau after step 10,000 suggests continued
improvement from lower learning rates and extended training.
Integrating reinforcement learning with a Romanian ASR reward signal
could directly optimise intelligibility beyond supervised fine-tuning.
Adapter-based hybrid approaches could combine fidelity preservation
of frozen modules with the phoneme-mapping strategy demonstrated here.
Finally, subjective evaluation with native Romanian listeners and
extension to spontaneous speech corpora remain important steps toward
a production-grade Romanian TTS system.

% ============================================================
\balance
\begin{thebibliography}{12}

\bibitem{valle}
C.~Wang, S.~Chen, Y.~Wu, Z.~Zhang, L.~Zhou, S.~Liu, Z.~Chen,
Y.~Liu, H.~Wang, J.~Li, L.~He, S.~Zhao, and F.~Wei,
``Neural Codec Language Models are Zero-Shot Text to Speech Synthesizers,''
\textit{arXiv preprint arXiv:2301.02111}, 2023.

\bibitem{seedtts}
P.~Anastassiou, J.~Chen, J.~Deng, Z.~Gao, Y.~Gu, Y.~He,
\textit{et al.},
``Seed-TTS: A Family of High-Quality Versatile Speech Generation Models,''
\textit{arXiv preprint arXiv:2406.02430}, 2024.

\bibitem{f5tts}
Y.~Chen, Z.~Ye, J.~Li, Z.~Huang, J.~Yang, S.~Li, T.~Bian,
B.~Wang, L.~Wu, Z.~Zhao, W.~Chen, and B.~Tian,
``F5-TTS: A Fairytaler that Fakes Fluent and Faithful Speech with Flow Matching,''
in \textit{Proc. ACL}, pp.\ 5904--5918, 2025.

\bibitem{cosyvoice}
Z.~Du, J.~Chen, X.~Guo, K.~Chen, X.~Zheng, Y.~Zheng, \textit{et al.},
``CosyVoice: A Scalable Multilingual Zero-shot Text-to-speech Synthesizer
based on Supervised Semantic Tokens,''
\textit{arXiv preprint arXiv:2407.05407}, 2024.

\bibitem{chatterbox}
Resemble AI,
``Chatterbox: State-of-the-Art Open Source TTS,''
GitHub, 2025. [Online]. Available:
\url{https://github.com/resemble-ai/chatterbox}

\bibitem{fastpitch_ro}
R.~R\u{a}gman and A.~Stan,
``Efficient Training Strategies for Natural Sounding Speech Synthesis
in Romanian,''
\textit{arXiv preprint arXiv:2410.06787}, 2024.

\bibitem{f5tts_ro}
A.~Chivereanu and T.~Boros,
``Extending F5-TTS to Romanian TTS via Lightweight Input Adaptation,''
\textit{arXiv preprint arXiv:2512.12297}, 2025.

\bibitem{mms}
V.~Pratap, A.~Tjandra, B.~Shi, P.~Tomasello, A.~Babu,
S.~Kundu, \textit{et al.},
``Scaling Speech Technology to 1,000+ Languages,''
\textit{J. Mach. Learn. Res.}, vol.~25, pp.~1--52, 2024.

\bibitem{align2speak}
J.~Chodroff, T.~Mahajan, V.~Pliakas, L.~Saul, and A.~Bhatt,
``Align2Speak: Improving TTS for Low Resource Languages via
ASR Preference Alignment,''
\textit{arXiv preprint arXiv:2509.21718}, 2025.

\bibitem{swara}
A.~Stan, J.~Yamagishi, S.~King, and M.~Aylett,
``The SWARA Speech Corpus: A Large Parallel Romanian Read Speech
Dataset,''
in \textit{Proc. SPED}, 2017.

\bibitem{utmos}
T.~Saeki, D.~Xin, W.~Nakata, T.~Koriyama, S.~Takamichi,
and H.~Saruwatari,
``UTMOS: UTokyo-SaruLab System for VoiceMOS Challenge 2022,''
in \textit{Proc. Interspeech}, pp.\ 4521--4525, 2022.

\bibitem{whisper}
A.~Radford, J.~W.~Kim, T.~Xu, G.~Brockman, C.~McLeavey,
and I.~Sutskever,
``Robust Speech Recognition via Large-Scale Weak Supervision,''
in \textit{Proc. ICML}, pp.\ 28492--28518, 2023.

\end{thebibliography}

\end{document}
